\documentclass[10pt,a4paper]{amsart}
\usepackage{amsmath,amssymb,amsthm,graphicx}
\usepackage{verbatim}
\newenvironment{reference}{\comment}{\endcomment}
%\newenvironment{reference}{}{}
\newenvironment{slogan}{\comment}{\endcomment}
\newenvironment{history}{\comment}{\endcomment}

% For commutative diagrams we use Xy-pic
\usepackage[all]{xy}

% We use 2cell for 2-commutative diagrams.
\xyoption{2cell}
% Theorem environments.
%
\theoremstyle{plain}
\newtheorem{theorem}[subsection]{Theorem}
\newtheorem{proposition}[subsection]{Proposition}
\newtheorem{lemma}[subsection]{Lemma}

\theoremstyle{definition}
\newtheorem{definition}[subsection]{Definition}
\newtheorem{example}[subsection]{Example}
\newtheorem{exercise}[subsection]{Exercise}
\newtheorem{situation}[subsection]{Situation}

\theoremstyle{remark}
\newtheorem{remark}[subsection]{Remark}
\newtheorem{remarks}[subsection]{Remarks}

\numberwithin{equation}{subsection}

% Macros
%
\def\lim{\mathop{\mathrm{lim}}\nolimits}
\def\colim{\mathop{\mathrm{colim}}\nolimits}
\def\Spec{\mathop{\mathrm{Spec}}}
\def\Hom{\mathop{\mathrm{Hom}}\nolimits}
\def\Ext{\mathop{\mathrm{Ext}}\nolimits}
\def\SheafHom{\mathop{\mathcal{H}\!\mathit{om}}\nolimits}
\def\SheafExt{\mathop{\mathcal{E}\!\mathit{xt}}\nolimits}
\def\Sch{\mathit{Sch}}
\def\Mor{\mathop{\mathrm{Mor}}\nolimits}
\def\Ob{\mathop{\mathrm{Ob}}\nolimits}
\def\Sh{\mathop{\mathit{Sh}}\nolimits}
\def\NL{\mathop{N\!L}\nolimits}
\def\CH{\mathop{\mathrm{CH}}\nolimits}
\def\proetale{{pro\text{-}\acute{e}tale}}
\def\etale{{\acute{e}tale}}
\def\QCoh{\mathit{QCoh}}
\def\Ker{\mathop{\mathrm{Ker}}}
\def\Im{\mathop{\mathrm{Im}}}
\def\Coker{\mathop{\mathrm{Coker}}}
\def\Coim{\mathop{\mathrm{Coim}}}

% Boxtimes
%
\DeclareMathSymbol{\boxtimes}{\mathbin}{AMSa}{"02}

%
% Macros for moduli stacks/spaces
%
\def\QCohstack{\mathcal{QC}\!\mathit{oh}}
\def\Cohstack{\mathcal{C}\!\mathit{oh}}
\def\Spacesstack{\mathcal{S}\!\mathit{paces}}
\def\Quotfunctor{\mathrm{Quot}}
\def\Hilbfunctor{\mathrm{Hilb}}
\def\Curvesstack{\mathcal{C}\!\mathit{urves}}
\def\Polarizedstack{\mathcal{P}\!\mathit{olarized}}
\def\Complexesstack{\mathcal{C}\!\mathit{omplexes}}
% \Pic is the operator that assigns to X its picard group, usage \Pic(X)
% \Picardstack_{X/B} denotes the Picard stack of X over B
% \Picardfunctor_{X/B} denotes the Picard functor of X over B
\def\Pic{\mathop{\mathrm{Pic}}\nolimits}
\def\Picardstack{\mathcal{P}\!\mathit{ic}}
\def\Picardfunctor{\mathrm{Pic}}
\def\Deformationcategory{\mathcal{D}\!\mathit{ef}}
\usepackage{cleveref}
\crefname{originallemma}{original lemma}{original lemmas}
\Crefname{originallemma}{Original lemma}{Original lemmas}
\makeatletter
\newlabel{morphisms-lemma-flat-base-change-scheme-theoretic-image}{{081I}{source}{}{sourcepointer-1}{}}
\newlabel{morphisms-lemma-flat-base-change-scheme-theoretic-image@cref}{{[originallemma][0][]081I}{[1][0][]source}{}{}{}}
\makeatother
\crefname{originallemma}{original lemma}{original lemmas}
\Crefname{originallemma}{Original lemma}{Original lemmas}
\makeatletter
\newlabel{morphisms-lemma-quasi-compact-immersion}{{01RG}{source}{}{sourcepointer-2}{}}
\newlabel{morphisms-lemma-quasi-compact-immersion@cref}{{[originallemma][0][]01RG}{[1][0][]source}{}{}{}}
\makeatother
\crefname{originallemma}{original lemma}{original lemmas}
\Crefname{originallemma}{Original lemma}{Original lemmas}
\makeatletter
\newlabel{spaces-properties-lemma-subspaces-presentation}{{07TW}{source}{}{sourcepointer-3}{}}
\newlabel{spaces-properties-lemma-subspaces-presentation@cref}{{[originallemma][0][]07TW}{[1][0][]source}{}{}{}}
\makeatother
\title{A quasi-compact immersion of algebraic spaces factors as an open immersion followed by a closed immersion}
\author{The Stacks Project Authors}
\date{stacks-project:spaces-morphisms.tex:local-snapshot-20261012; GNU FDL 1.2 or later}
\begin{document}
\maketitle
\setcounter{subsection}{1}
\noindent\textit{Editorial note.} All delivered mathematical text is the author's own, copied from the pinned Stacks Project chapter file spaces-morphisms.tex; no statement or proof step was rewritten, reordered or removed.\par\medskip
\noindent\textit{Editorial note.} The proof is a descent along a presentation; the flat base change property of scheme theoretic images is the step that makes the descent possible.\par\medskip
\noindent\textit{Editorial note.} The three references of the counted spans are recorded as non-load-bearing navigation pointers; the argument that consumes them stays inside the counted proof span.\par\medskip
\noindent\textit{Editorial note.} Preamble: the three native excerpt spans of sources/stacks-spaces-morphisms/preamble.tex (lines 8-13, 14-18 and 154-217). The chapter file itself carries no preamble text.\par\medskip
\noindent\textit{Allowed standard prerequisites (editorial).}\begin{itemize}
\item Scheme theoretic images of quasi-compact morphisms commute with flat base change. Morphisms, Lemma morphisms-lemma-flat-base-change-scheme-theoretic-image, used at proof lines 2386-2391. Gives the equality s\textasciicircum{}\{-1\}Tbar = t\textasciicircum{}\{-1\}Tbar needed for descent.
\item Descending subspaces through a presentation of an algebraic space. Properties of Spaces, Lemma spaces-properties-lemma-subspaces-presentation, used at proof lines 2392-2395. Produces the closed subspace Zbar of X with pullback Tbar.
\item The scheme case: a quasi-compact immersion is an open immersion onto its scheme theoretic image. Morphisms, Lemma morphisms-lemma-quasi-compact-immersion, used at proof lines 2396-2398. Makes T -> Tbar an open immersion, hence Z -> Zbar as well.
\item Surjective etale presentations of algebraic spaces. used at proof lines 2384-2386. Supplies U, R and the projections used throughout.
\end{itemize}\medskip
\ifdefined\hypertarget\hypertarget{sourcepointer-1}{}\fi
\noindent\textit{Source reference.} Original lemma 081I: focused-seeds/AGP/sources/stacks-snapshot-20261011/tags.txt, source lines 10218--10218. Navigation pointer into the Stacks Project; the standard fact it supplies is applied inside the counted proof span and the argument resting on it stays there.\par\medskip
\ifdefined\hypertarget\hypertarget{sourcepointer-2}{}\fi
\noindent\textit{Source reference.} Original lemma 01RG: focused-seeds/AGP/sources/stacks-snapshot-20261011/tags.txt, source lines 2301--2301. Navigation pointer into the Stacks Project; the standard fact it supplies is applied inside the counted proof span and the argument resting on it stays there.\par\medskip
\ifdefined\hypertarget\hypertarget{sourcepointer-3}{}\fi
\noindent\textit{Source reference.} Original lemma 07TW: focused-seeds/AGP/sources/stacks-snapshot-20261011/tags.txt, source lines 9951--9951. Navigation pointer into the Stacks Project; the standard fact it supplies is applied inside the counted proof span and the argument resting on it stays there.\par\medskip
\begin{lemma}
\label{lemma-factor-the-other-way}
Let $S$ be a scheme. Let $Z \to X$ be an immersion of algebraic spaces over
$S$. Assume $Z \to X$ is quasi-compact.
There exists a factorization $Z \to \overline{Z} \to X$ where
$Z \to \overline{Z}$ is an open immersion and $\overline{Z} \to X$
is a closed immersion.
\end{lemma}

\begin{proof}
Let $U$ be a scheme and let $U \to X$ be surjective \'etale.
As usual denote $R = U \times_X U$ with projections
$s, t : R \to U$. Set $T = Z \times_U X$. Let $\overline{T} \subset U$
be the scheme theoretic image of $T \to U$. Note that
$s^{-1}\overline{T} = t^{-1}\overline{T}$ as taking
scheme theoretic images of quasi-compact morphisms commute
with flat base change, see
Morphisms, Lemma \ref{morphisms-lemma-flat-base-change-scheme-theoretic-image}.
Hence we obtain a closed subspace $\overline{Z} \subset X$ whose
pullback to $U$ is $\overline{T}$, see
Properties of Spaces, Lemma
\ref{spaces-properties-lemma-subspaces-presentation}.
By Morphisms, Lemma \ref{morphisms-lemma-quasi-compact-immersion}
the morphism $T \to \overline{T}$
is an open immersion. It follows that $Z \to \overline{Z}$ is
an open immersion and we win.
\end{proof}

\end{document}
